% Referencias. Las marcadas \pv requieren comprobación antes de la entrega (contrastar cada DOI en Crossref).
\begin{thebibliography}{99}
\addcontentsline{toc}{chapter}{Referencias}

\bibitem[Bennett(2026)]{bennett2026}
Bennett, D. (2026). \emph{How Angine de Poitrine use Microtonality} [video]. David Bennett Music Theory, YouTube,
21 de marzo de 2026. \url{https://www.youtube.com/watch?v=91QWp6Ygx7c}

\bibitem[Bittner et~al.(2022)]{bittner2022}
Bittner, R. M., Bosch, J. J., Rubinstein, D., Meseguer-Brocal, G., y Ewert, S. (2022). A lightweight
instrument-agnostic model for polyphonic note transcription and multipitch estimation. \emph{Proc.\ IEEE ICASSP},
781--785. \url{https://doi.org/10.1109/ICASSP43922.2022.9746549}

\bibitem[Bozkurt(2008)]{bozkurt2008}
Bozkurt, B. (2008). An automatic pitch analysis method for Turkish maqam music. \emph{Journal of New Music Research},
37(1), 1--13.

\bibitem[Bozkurt et~al.(2009)]{bozkurt2009}
Bozkurt, B., Yarman, O., Karaosmanoğlu, M. K., y Akkoç, C. (2009). Weighing diverse theoretical models on Turkish maqam
music against pitch measurements. \emph{Journal of New Music Research}, 38(1), 45--70.
\url{https://doi.org/10.1080/09298210903147673}

\bibitem[Bozkurt et~al.(2014)]{bozkurt2014}
Bozkurt, B., Ayangil, R., y Holzapfel, A. (2014). Computational analysis of Turkish makam music: Review of
state-of-the-art and challenges. \emph{Journal of New Music Research}, 43(1), 3--23.

\bibitem[Broze y Shanahan(2013)]{broze2013}
Broze, Y., y Shanahan, D. (2013). Diachronic changes in jazz harmony: A cognitive perspective. \emph{Music Perception},
31(1), 32--45. \url{https://doi.org/10.1525/mp.2013.31.1.32}

\bibitem[Chorościan(2018)]{choroscian2018}
Chorościan, E. (2018). Alois Hába's microtonal systems and their practical use in selected string quartets.
\emph{Kwartalnik Młodych Muzykologów UJ}, 38(3), 47--78. \url{https://doi.org/10.4467/23537094KMMUJ.18.033.9734}

\bibitem[Davis y Mermelstein(1980)]{davis1980}
Davis, S. B., y Mermelstein, P. (1980). Comparison of parametric representations for monosyllabic word recognition in
continuously spoken sentences. \emph{IEEE Transactions on Acoustics, Speech, and Signal Processing}, 28(4), 357--366.
\url{https://doi.org/10.1109/TASSP.1980.1163420}

\bibitem[de Cheveigné y Kawahara(2002)]{decheveigne2002}
de Cheveigné, A., y Kawahara, H. (2002). YIN, a fundamental frequency estimator for speech and music. \emph{Journal of
the Acoustical Society of America}, 111(4), 1917--1930. \url{https://doi.org/10.1121/1.1458024}

\bibitem[de Clercq y Temperley(2011)]{declercq2011}
de Clercq, T., y Temperley, D. (2011). A corpus analysis of rock harmony. \emph{Popular Music}, 30(1), 47--70.
\url{https://doi.org/10.1017/S026114301000067X}

\bibitem[Dunlevie(2026)]{dunlevie2026}
Dunlevie, J. (2026). Angine de Poitrine, the French Canadian costumed microtonal duo breaking the internet. \emph{ABC
News}, 20 de mayo de 2026.
\url{https://www.abc.net.au/news/2026-05-20/angine-de-poitrine-microtonal-math-rock-quebec/106593634}

\bibitem[Ellis(2007)]{ellis2007}
Ellis, D. P. W. (2007). Beat tracking by dynamic programming. \emph{Journal of New Music Research}, 36(1), 51--60.
\url{https://doi.org/10.1080/09298210701653344}

\bibitem[Endres y Schindelin(2003)]{endres2003}
Endres, D. M., y Schindelin, J. E. (2003). A new metric for probability distributions. \emph{IEEE Transactions on
Information Theory}, 49(7), 1858--1860. \url{https://doi.org/10.1109/TIT.2003.813506}

\bibitem[Farhat(1990)]{farhat1990}
Farhat, H. (1990). \emph{The Dastgah Concept in Persian Music}. Cambridge University Press.
\url{https://doi.org/10.1017/CBO9780511470233}

\bibitem[Gedik y Bozkurt(2010)]{gedik2010}
Gedik, A. C., y Bozkurt, B. (2010). Pitch-frequency histogram-based music information retrieval for Turkish music.
\emph{Signal Processing}, 90(4), 1049--1063. \url{https://doi.org/10.1016/j.sigpro.2009.06.017}

\bibitem[Hába(1927)]{haba1927}
Hába, A. (1927). \emph{Neue Harmonielehre des diatonischen, chromatischen Viertel-, Drittel-, Sechstel- und
Zwölftel-Tonsystems}. Traducción del checo por el autor, revisión de E. Steinhard. Fr.\ Kistner \& C.\ F.\ W. Siegel.

\bibitem[Holzapfel et~al.(2014)]{holzapfel2014}
Holzapfel, A., Krebs, F., y Srinivasamurthy, A. (2014). Tracking the «odd»: Meter inference in a culturally diverse
music corpus. \emph{Proc.\ ISMIR}, 425--430.

\bibitem[Hooktheory(s.\,f.)]{hooktheory}
Hooktheory (s.\,f.). \emph{TheoryTab}: análisis armónicos de los álbumes \emph{30} (Adele, 2021), \emph{Everything to
Everyone} (Reneé Rapp, 2022), \emph{Gettin' Old} (Luke Combs, 2023), \emph{Happier Than Ever} (Billie Eilish, 2021),
\emph{Midnights} (Taylor Swift, 2022), \emph{Montero} (Lil Nas X, 2021), \emph{Solar Power} (Lorde, 2021),
\emph{$-$ (Subtract)} (Ed Sheeran, 2023) y \emph{Think Later} (Tate McRae, 2023). \url{https://www.hooktheory.com/theorytab}
\pv[: enlace exacto de cada análisis de álbum]

\bibitem[Huron(2006)]{huron2006}
Huron, D. (2006). \emph{Sweet Anticipation: Music and the Psychology of Expectation}. MIT Press.

\bibitem[Jin(2015)]{jin2015}
Jin, Y. (2015). \emph{A Study of Quarter-Tone Music for Solo Violin by Alois Hába} [tesis de doctorado en artes
musicales]. College-Conservatory of Music, University of Cincinnati.

\bibitem[Kanani et~al.(2025)]{kanani2025}
Kanani, M., O'Leary, S., y McDermott, J. (2025). Radif Corpus: A symbolic dataset for non-metric Iranian classical
music. \emph{Proc.\ ISMIR}. arXiv:2507.10456.

\bibitem[KEXP(2026)]{kexp2026}
KEXP (2026). \emph{Angine de Poitrine: KEXP session feature} [artículo], 25 de marzo de 2026.
\url{https://www.kexp.org/read/2026/3/25/angine-de-poitrine-kexp-session-feature/}

\bibitem[Karaosmanoğlu(2012)]{karaosmanoglu2012}
Karaosmanoğlu, M. K. (2012). A Turkish makam music symbolic database for music information retrieval: SymbTr.
\emph{Proc.\ ISMIR}, 223--228.

\bibitem[Kim et~al.(2018)]{kim2018}
Kim, J. W., Salamon, J., Li, P., y Bello, J. P. (2018). CREPE: A convolutional representation for pitch estimation.
\emph{Proc.\ IEEE ICASSP}. \url{https://doi.org/10.1109/ICASSP.2018.8461329}

\bibitem[King Gizzard and the Lizard Wizard(2017)]{kinggizzard2017}
King Gizzard and the Lizard Wizard (2017). \emph{Flying Microtonal Banana} [álbum y nota de lanzamiento], 24 de febrero
de 2017. \url{https://kinggizzardandthelizardwizard.com/release/flying-microtonal-banana}

\bibitem[Krumhansl(1990)]{krumhansl1990}
Krumhansl, C. L. (1990). \emph{Cognitive Foundations of Musical Pitch}. Oxford University Press.

\bibitem[Lendvai(1971)]{lendvai1971}
Lendvai, E. (1971). \emph{Béla Bartók: An Analysis of His Music}. Kahn \& Averill.


\bibitem[Lerdahl y Jackendoff(1983)]{lerdahl1983}
Lerdahl, F., y Jackendoff, R. (1983). \emph{A Generative Theory of Tonal Music}. MIT Press.

\bibitem[Lin(1991)]{lin1991}
Lin, J. (1991). Divergence measures based on the Shannon entropy. \emph{IEEE Transactions on Information Theory},
37(1), 145--151. \url{https://doi.org/10.1109/18.61115}

\bibitem[Little(2019)]{little2019}
Little, T. (2019). \emph{The Quarter-Tone Problem (and Hába's «Field Shifting» Solution)} [video]. Classical Nerd,
YouTube, 17 de enero de 2019. \url{https://www.youtube.com/watch?v=E_Z0q7jqyqI}

\bibitem[Loth et~al.(2023)]{loth2023}
Loth, J., Sarmento, P., Carr, C. J., Zukowski, Z., y Barthet, M. (2023). ProgGP: From GuitarPro tablature neural
generation to progressive metal production. \emph{Proc.\ CMMR}. arXiv:2307.05328.

\bibitem[Lu et~al.(2024)]{lu2024}
Lu, W.-T., Wang, J.-C., Kong, Q., y Hung, Y.-N. (2024). Music source separation with band-split RoPE transformer.
\emph{Proc.\ IEEE ICASSP}, 481--485. \url{https://doi.org/10.1109/ICASSP48485.2024.10446843}

\bibitem[Mathémusique(2026)]{mathemusique2026}
Mathémusique [Lascabettes, P.] (2026). \emph{Comprendre ANGINE DE POITRINE} [video]. YouTube, 16 de mayo de 2026.
\url{https://www.youtube.com/watch?v=96D7DKZh-es}

\bibitem[Mauch y Dixon(2014)]{mauch2014}
Mauch, M., y Dixon, S. (2014). pYIN: A fundamental frequency estimator using probabilistic threshold distributions.
\emph{Proc.\ IEEE ICASSP}, 659--663. \url{https://doi.org/10.1109/ICASSP.2014.6853678}

\bibitem[McClain(2026)]{mcclain2026}
McClain, L. (2026). \emph{Angine de Poitrine Explained} [video]. YouTube, 7 de septiembre de 2026.
\url{https://www.youtube.com/watch?v=kWv5zjsjoHw}

\bibitem[McFee et~al.(2015)]{mcfee2015}
McFee, B., Raffel, C., Liang, D., Ellis, D. P. W., McVicar, M., Battenberg, E., y Nieto, O. (2015). librosa: Audio and
music signal analysis in Python. \emph{Proc.\ SciPy}, 18--24. \url{https://doi.org/10.25080/Majora-7b98e3ed-003}

\bibitem[MIDI Manufacturers Association(1996)]{midi1996}
MIDI Manufacturers Association (1996). \emph{The Complete MIDI 1.0 Detailed Specification} (versión 96.1). MMA.

\bibitem[Monty y Benabbou(s.\,f.)]{noize}
Monty, A., y Benabbou, G. (s.\,f.). Angine De Poitrine: A discussion with the infectious Quebec band. \emph{Noize
Magazine}, 20 de febrero (la página no indica el año; por el contexto, 2026).
\url{https://www.noizemagazine.com/blog/angine-de-poitrine-a-discussion-with-the-infectious-quebec-band}

\bibitem[Khn de Poitrine(2026)]{aucard2026}
Khn de Poitrine (2026). [Entrevista] Angine de Poitrine. \emph{Aucard de Tours} [podcast], 26 de febrero de 2026,
1:33--2:58. \url{https://podcastaddict.com/aucard-de-tours/episode/218290057}

\bibitem[NIST(2015)]{nist2015}
National Institute of Standards and Technology (2015). \emph{Secure Hash Standard (SHS)}. FIPS PUB 180-4.
\url{https://doi.org/10.6028/NIST.FIPS.180-4}

\bibitem[Pearce(2018)]{pearce2018}
Pearce, M. T. (2018). Statistical learning and probabilistic prediction in music cognition: Mechanisms of stylistic
enculturation. \emph{Annals of the New York Academy of Sciences}, 1423(1), 378--395.
\url{https://doi.org/10.1111/nyas.13654}

\bibitem[Piston(1987)]{piston1987}
Piston, W. (1987). \emph{Harmony} (5.ª ed., revisada y ampliada por M. DeVoto). W. W. Norton.

\bibitem[Rafii y Pardo(2013)]{rafii2013}
Rafii, Z., y Pardo, B. (2013). REpeating Pattern Extraction Technique (REPET): A simple method for music/voice
separation. \emph{IEEE Transactions on Audio, Speech, and Language Processing}, 21(1), 71--82.
\url{https://doi.org/10.1109/TASL.2012.2213249}

\bibitem[Riemann(1893)]{riemann1893}
Riemann, H. (1893). \emph{Vereinfachte Harmonielehre oder die Lehre von den tonalen Funktionen der Akkorde}. Augener, Londres.

\bibitem[Rouard et~al.(2023)]{rouard2023}
Rouard, S., Massa, F., y Défossez, A. (2023). Hybrid transformers for music source separation. \emph{Proc.\ IEEE
ICASSP}. \url{https://doi.org/10.1109/ICASSP49357.2023.10096956}

\bibitem[Salamon y Gómez(2012)]{salamon2012}
Salamon, J., y Gómez, E. (2012). Melody extraction from polyphonic music signals using pitch contour characteristics.
\emph{IEEE Transactions on Audio, Speech, and Language Processing}, 20(6), 1759--1770.
\url{https://doi.org/10.1109/TASL.2012.2188515}

\bibitem[Serra(2011)]{serra2011}
Serra, X. (2011). A multicultural approach in music information research. \emph{Proc.\ 12th International Society for
Music Information Retrieval Conference (ISMIR)}, Miami, 151--156.

\bibitem[Shanahan et~al.(2012)]{shanahan2012}
Shanahan, D., Broze, Y., y Rodgers, R. (2012). A diachronic analysis of harmonic schemata in jazz. \emph{Proc.\ ICMPC 12
/ ESCOM 8}, 909--917.

\bibitem[Shanahan y Broze(s.\,f.)]{irealpro}
Shanahan, D., y Broze, Y. (s.\,f.). \emph{iRealPro Corpus of Jazz Standards} (v1.0). Zenodo.
\url{https://doi.org/10.5281/zenodo.3546040}

\bibitem[Skinner(2006)]{skinner2006}
Skinner, M. L. (2006). \emph{Toward a Quarter-Tone Syntax: Analyses of Selected Works by Blackwood, Hába, Ives, and
Wyschnegradsky} [tesis doctoral]. State University of New York at Buffalo.

\bibitem[Spurný(2021)]{spurny2021}
Spurný, L. (2021). Six observations on the performance of Hába's microtonal music. En H. Spurná y K. St.~Pierre (eds.),
\emph{New Paths in Opera: Martinů, Burian, Hába, Schulhoff, Ullmann} (pp.~221--240). KLP / Hollitzer.

\bibitem[Singh et~al.(2022)]{singh2022}
Singh, Y., Gupta, Y., Patar, S., Saraswat, A., y Biswas, A. (2022). \emph{Hindustani Classical Music Transcription
Dataset} (v1.0.1). Zenodo. \url{https://doi.org/10.5281/zenodo.7424861}

\bibitem[Srinivasamurthy et~al.(2016)]{srinivasamurthy2016}
Srinivasamurthy, A., Holzapfel, A., Cemgil, A. T., y Serra, X. (2016). A generalized Bayesian model for tracking long
metrical cycles in acoustic music signals. \emph{Proc.\ IEEE ICASSP}, 76--80.
\url{https://doi.org/10.1109/ICASSP.2016.7471640}

\bibitem[Stevens et~al.(1937)]{stevens1937}
Stevens, S. S., Volkmann, J., y Newman, E. B. (1937). A scale for the measurement of the psychological magnitude pitch.
\emph{Journal of the Acoustical Society of America}, 8(3), 185--190. \url{https://doi.org/10.1121/1.1915893}

\bibitem[Temperley(2009)]{temperley2009}
Temperley, D. (2009). \emph{A Statistical Analysis of Tonal Harmony} [página web con el corpus Kostka-Payne: 46
fragmentos de Kostka y Payne, \emph{Tonal Harmony}, 3.ª ed., McGraw-Hill, 1995]. \url{https://davidtemperley.com/kp-stats/}

\bibitem[Temperley y de Clercq(2013)]{temperley2013}
Temperley, D., y de Clercq, T. (2013). Statistical analysis of harmony and melody in rock music. \emph{Journal of New
Music Research}, 42(3), 187--204.

\bibitem[Touma(1996)]{touma1996}
Touma, H. H. (1996). \emph{The Music of the Arabs} (trad. L. Schwartz, ed. ampliada). Amadeus Press.

\bibitem[Weller(2026)]{weller2026}
Weller, P. (2026). «The moment we started playing with it, we just laughed»: Viral duo Angine de Poitrine have become one
of the most talked-about bands on the planet -- and microtonal guitar has been the key to their unlikely success.
\emph{Guitar World}, 20 de marzo de 2026. \url{https://www.guitarworld.com/artists/bands/angine-de-poitrine-microtonal-guitar}

\bibitem[van Rijsbergen(1979)]{vanrijsbergen1979}
van Rijsbergen, C. J. (1979). \emph{Information Retrieval} (2.ª ed.). Butterworths.

\bibitem[Viterbi(1967)]{viterbi1967}
Viterbi, A. J. (1967). Error bounds for convolutional codes and an asymptotically optimum decoding algorithm.
\emph{IEEE Transactions on Information Theory}, 13(2), 260--269. \url{https://doi.org/10.1109/TIT.1967.1054010}

\end{thebibliography}

